\chapter{مدیریت و پیکربندی شبکه}

در قلمرو شبکه، عملاً هیچ بن‌بستی برای لینوکس وجود ندارد؛ این سیستم‌عامل زیربنای طیف وسیعی از تجهیزات و زیرساخت‌های شبکه‌ای—از دیوارهای آتش (\en{Firewalls}) و مسیریاب‌ها (\en{Routers}) گرفته تا سرویس‌دهنده‌های نام دامنه (\en{DNS Servers}) و ذخیره‌سازهای متصل به شبکه (\en{Network-Attached Storage - NAS})—را تشکیل داده است.

با توجه به گستردگی مفاهیم شبکه، ابزارهای متعددی نیز برای پیکربندی و پایش آن طراحی شده‌اند. تمرکز ما در این بخش بر پرکاربردترین ابزارهایی است که برای پایش ترافیک، تشخیص عیوب اتصال، انتقال فایل، و مدیریت از راه دور استفاده می‌شوند. در این فصل، ابزارهای زیر را بررسی خواهیم کرد:

\begin{itemize}
  \item \cmd{ping}: ارسال بسته‌های \en{ICMP} از نوع \en{ECHO\_REQUEST} به یک میزبان (\en{Host}) مشخص، به‌منظور سنجش در دسترس‌بودن آن و اندازه‌گیری زمان رفت‌وبرگشت بسته (\en{Round-Trip Time})؛ ساده‌ترین و بنیادی‌ترین ابزار برای عیب‌یابی اولیه‌ی اتصال شبکه.
  \item \cmd{traceroute}: رهگیری و نمایشِ گام‌به‌گام مسیر پیموده‌شده توسط بسته‌ها—شامل تمامی مسیریاب‌های میانی (\en{Hops})—تا رسیدن به میزبان مقصد؛ ابزاری کلیدی برای شناسایی محل دقیق تأخیر یا قطعی ارتباط در طول مسیر.
  \item \cmd{ip}: ابزاری جامع و قدرتمند برای نمایش و مدیریتِ جدول‌های مسیریابی (\en{Routing Tables})، رابط‌های شبکه (\en{Network Interfaces})، آدرس‌های \en{IP}، و تونل‌ها؛ که در توزیع‌های مدرن، عملاً جایگزین ابزارهای قدیمی‌تری مانند \cmd{ifconfig} و \cmd{route} شده است.
  \item \cmd{netstat}: پایش سوکت‌های شبکه (\en{Network Sockets})، جدول‌های مسیریابی، آمار رابط‌های شبکه، و وضعیت اتصالات فعال (نظیر اتصالات در حال گوش‌دادن یا برقرار)؛ ابزاری مفید برای تشخیص اینکه چه سرویس‌ها و پورت‌هایی در حال استفاده هستند.
  \item \cmd{ftp}: کلاینت استاندارد پروتکل انتقال فایل (\en{File Transfer Protocol}) در شبکه، برای بارگذاری و دریافت فایل از سرورهای \en{FTP}.
  \item \cmd{curl}: ابزاری چندمنظوره و خط‌فرمانی برای انتقال داده از طریق نشانی‌های \en{URL}، با پشتیبانی از طیف گسترده‌ای از پروتکل‌ها (نظیر \en{HTTP}، \en{HTTPS}، و \en{FTP})؛ پرکاربرد در اسکریپت‌نویسی و آزمایش \en{API}ها.
  \item \cmd{wget}: ابزار دانلود شبکه‌ای، با قابلیت اجرا در پس‌زمینه (\en{Background}) و پشتیبانی از حالتِ غیرتعاملی (\en{Non-interactive})؛ مناسب برای دانلود فایل‌های حجیم یا اسکریپت‌های خودکار.
  \item \cmd{ssh}: کلاینت پروتکل امن \en{OpenSSH}، برای برقراری اتصال رمزنگاری‌شده و دسترسی از راه دور به یک میزبان دیگر.
\end{itemize}

پیش‌فرض ما بر این است که شما با مبانی شبکه آشنایی دارید. در عصر حاضر، درک مفاهیم پایه‌ای زیر برای هر کاربر رایانه ضروری است:

\begin{itemize}
  \item نشانی پروتکل اینترنت (\en{IP Address})
  \item نام میزبان (\en{Hostname}) و نام دامنه (\en{Domain})
  \item شناسه یکنواخت منبع (\en{Uniform Resource Identifier - URI})
\end{itemize}

جهت تعمیق دانش خود در این زمینه‌ها، می‌توانید به بخش «منابع تکمیلی» در انتهای این فصل مراجعه کنید.

\begin{note}
\textbf{نکته:} برخی از ابزارهای مذکور ممکن است به‌صورت پیش‌فرض در توزیع شما نصب نباشند و نیاز به دریافت از مخازن رسمی داشته باشند. همچنین، اجرای تعدادی از این فرامین مستلزم برخورداری از امتیازات ابرکاربر (\en{Superuser}) است.
\end{note}

\section{ارزیابی و پایش شبکه}

حتی اگر مدیر سیستم نباشید، اغلب مفید است که عملکرد و وضعیت یک شبکه را بررسی کنید.

\subsection{عیب‌یابی اولیه با ابزار ping}

نخستین و بنیادی‌ترین ابزار در تحلیل شبکه، دستور \cmd{ping} است. این فرمان یک بسته‌ی ویژه‌ی تحت پروتکل \en{ICMP} با عنوان \en{ECHO\_REQUEST} را به مقصد یک میزبان (\en{Host}) مشخص ارسال می‌کند. اغلب دستگاه‌های شبکه پس از دریافت این بسته، با یک پاسخ (\en{Reply}) واکنش نشان می‌دهند که همین فرآیند ساده، امکان سنجش پایداری اتصال را فراهم می‌آورد.

\begin{note}
\textbf{نکته:} شایان ذکر است که بسیاری از مدیران شبکه، دستگاه‌ها (از جمله سرورهای لینوکس) را به‌گونه‌ای پیکربندی می‌کنند که این بسته‌ها را نادیده بگیرند. این رویکرد معمولاً با هدف مخفی‌سازی میزبان از دید مهاجمان یا از طریق محدودیت‌های اعمال شده در دیوارهای آتش (\en{Firewalls}) صورت می‌گیرد.
\end{note}

به‌عنوان مثال، برای ارزیابی دسترسی‌پذیری وب‌سایت \file{linuxcommand.org}، می‌توان دستور را به شکل زیر اجرا کرد:

\begin{latin}
\begin{lstlisting}
[me@linuxbox ~]$ ping linuxcommand.org
\end{lstlisting}
\end{latin}

پس از اجرا، \cmd{ping} به‌صورت مستمر و در بازه‌های زمانی معین (به‌طور پیش‌فرض هر یک ثانیه) اقدام به ارسال بسته می‌کند تا زمانی که توسط کاربر متوقف شود.

\begin{latin}
\begin{lstlisting}
[me@linuxbox ~]$ ping linuxcommand.org
PING linuxcommand.org (66.35.250.210) 56(84) bytes of data.
64 bytes from vhost.sourceforge.net (66.35.250.210): icmp_seq=1 ttl=43 time=107 ms
64 bytes from vhost.sourceforge.net (66.35.250.210): icmp_seq=2 ttl=43 time=108 ms
64 bytes from vhost.sourceforge.net (66.35.250.210): icmp_seq=3 ttl=43 time=106 ms
64 bytes from vhost.sourceforge.net (66.35.250.210): icmp_seq=4 ttl=43 time=106 ms
64 bytes from vhost.sourceforge.net (66.35.250.210): icmp_seq=5 ttl=43 time=105 ms
64 bytes from vhost.sourceforge.net (66.35.250.210): icmp_seq=6 ttl=43 time=107 ms

--- linuxcommand.org ping statistics ---
6 packets transmitted, 6 received, 0% packet loss, time 6010ms
rtt min/avg/max/mdev = 105.647/107.052/108.118/0.824 ms
\end{lstlisting}
\end{latin}

با فشردن کلید ترکیبی \cmd{Ctrl+C} عملیات متوقف شده و گزارش آماری عملکرد شبکه ارائه می‌گردد. در یک بستر شبکه‌ای ایده‌آل، نرخ از دست رفتن بسته‌ها (\en{Packet Loss}) باید صفر درصد باشد. دریافت پاسخ‌های موفقیت‌آمیز نشان‌دهنده‌ی سلامت کلی زیرساخت‌های فیزیکی و منطقی، شامل کارت‌های رابط شبکه (\en{Network Interface Controller - NIC})، کابل‌کشی، مسیریاب‌ها (\en{Routers}) و درگاه‌های خروجی (\en{Gateways}) است.

\subsection{ردیابی مسیر بسته‌ها با traceroute}

برنامه \cmd{traceroute} (و در برخی توزیع‌ها ابزار مشابه آن یعنی \cmd{tracepath}) گزارشی جامع از تمامی «گام‌ها» (\en{Hops}) یا مسیریاب‌هایی که ترافیک شبکه برای رسیدن از مبدأ به میزبان مقصد طی می‌کند، ارائه می‌دهد. به‌طور مثال، جهت واکاوی مسیر عبور بسته‌ها برای دسترسی به وب‌سایت \file{slashdot.org} دستور زیر را اجرا می‌کنیم:

\begin{latin}
\begin{lstlisting}
[me@linuxbox ~]$ traceroute slashdot.org
\end{lstlisting}
\end{latin}

خروجی حاصل بدین صورت خواهد بود:

\begin{latin}
\begin{lstlisting}
traceroute to slashdot.org (216.34.181.45), 30 hops max, 40 byte packets
 1  ipcop.localdomain (192.168.1.1)  1.066 ms  1.366 ms  1.720 ms
 2  * * *
 3  ge-4-13-ur01.rockville.md.bad.comcast.net (68.87.130.9)  14.622 ms  14.885 ms  15.169 ms
 4  po-30-ur02.rockville.md.bad.comcast.net (68.87.129.154)  17.634 ms  17.626 ms  17.899 ms
 5  po-60-or03.rockville.md.bad.comcast.net (68.87.129.158)  15.992 ms  15.983 ms  16.256 ms
 6  po-30-ar01.howardcounty.md.bad.comcast.net (68.87.136.5)  22.835 ms  14.233 ms  14.405 ms
 7  po-10-ar02.whitemarsh.md.bad.comcast.net (68.87.129.34)  16.154 ms  13.600 ms  18.867 ms
 8  te-0-3-0-1-cr01.philadelphia.pa.ibone.comcast.net (68.86.90.77)  21.951 ms  21.073 ms  21.557 ms
 9  pos-0-8-0-0-cr01.newyork.ny.ibone.comcast.net (68.86.85.10)  22.917 ms  21.884 ms  22.126 ms
10  204.70.144.1 (204.70.144.1)  43.110 ms  21.248 ms  21.264 ms
11  cr1-pos-0-7-3-1.newyork.savvis.net (204.70.195.93)  21.857 ms  cr2-pos-0-7-3-1.newyork.savvis.net (204.70.204.238)  19.556 ms  cr1-pos-0-7-3-1.newyork.savvis.net (204.70.195.93)  19.634 ms
12  cr2-pos-0-7-3-0.chicago.savvis.net (204.70.192.109)  41.586 ms  42.843 ms  cr2-tengig-0-0-2-0.chicago.savvis.net (204.70.196.242)  43.115 ms
13  hr2-tengigabitethernet-12-1.elkgrovech3.savvis.net (204.70.195.122)  44.215 ms  41.833 ms  45.658 ms
14  csr1-ve241.elkgrovech3.savvis.net (216.64.194.42)  46.840 ms  43.372 ms  47.041 ms
15  64.27.160.194 (64.27.160.194)  56.137 ms  55.887 ms  52.810 ms
16  slashdot.org (216.34.181.45)  42.727 ms  42.016 ms  41.437 ms
\end{lstlisting}
\end{latin}

مطابق این خروجی، اتصال میان سیستم مبدأ و مقصد نهایی مستلزم عبور از ۱۶ مسیریاب (\en{Router}) مختلف است. برای هر گره‌ای که اطلاعات هویتی خود را افشا می‌کند، نام میزبان (\en{Hostname})، نشانی \en{IP} و سه نمونه از زمان رفت‌وبرگشت بسته‌ها (\en{Round-Trip Time}) نمایش داده می‌شود. در مواردی که یک مسیریاب به دلایلی نظیر پیکربندی‌های امنیتی، ترافیک بالا یا محدودیت‌های دیوار آتش پاسخ ندهد، به‌جای اطلاعات فنی، نماد ستاره (\cmd{*}) درج می‌شود (مانند گام شماره ۲ در مثال فوق).

در شرایطی که اطلاعات مسیر توسط تجهیزات میانی مسدود شده باشد، گاهی می‌توان با استفاده از گزینه‌های \cmd{-T} (استفاده از پروتکل \en{TCP}) یا \cmd{-I} (استفاده از پروتکل \en{ICMP}) از این محدودیت‌ها عبور کرده و به اطلاعات مسیر دست یافت.

\subsection{مدیریت زیرساخت شبکه با ابزار ip}

فرمان \cmd{ip} یک ابزار پیکربندی شبکه چندمنظوره و پیشرفته است که از قابلیت‌های نوین هسته‌ی لینوکس (\en{Linux Kernel}) بهره می‌گیرد. این ابزار به‌طور کامل جایگزین ابزار کلاسیک و منسوخ‌شده‌ی \cmd{ifconfig} شده است. از برنامه \cmd{ip} جهت واکاوی تنظیمات، مشاهده‌ی آمار ترافیکی و مدیریت جنبه‌های مختلف معماری شبکه استفاده می‌شود. با بهره‌گیری از زیرفرمان‌های این ابزار، می‌توان وضعیت رابط‌های شبکه (\en{Interfaces}) و جداول مسیریابی (\en{Routing Tables}) سیستم را با دقت بالا تحلیل کرد.

در گام نخست، برای بررسی وضعیت رابط‌های شبکه از دستور زیر استفاده می‌کنیم، خروجی نمونه به شرح زیر است:

\begin{latin}
\begin{lstlisting}
[me@linuxbox ~]$ ip address show
1: lo: <LOOPBACK,UP,LOWER_UP> mtu 65536 qdisc noqueue state UNKNOWN group default qlen 1000
    link/loopback 00:00:00:00:00:00 brd 00:00:00:00:00:00
    inet 127.0.0.1/8 scope host lo
       valid_lft forever preferred_lft forever
    inet6 ::1/128 scope host noprefixroute
       valid_lft forever preferred_lft forever
2: enp1s0: <BROADCAST,MULTICAST,UP,LOWER_UP> mtu 1500 qdisc mq state UP group default qlen 1000
    link/ether 00:26:6c:26:67:bf brd ff:ff:ff:ff:ff:ff
    inet 192.168.1.223/24 brd 192.168.50.255 scope global dynamic noprefixroute enp1s0
       valid_lft 82366sec preferred_lft 82366sec
    inet6 fe80::226:6cff:fe26:67bf/64 scope link noprefixroute
       valid_lft forever preferred_lft forever
3: wlp2s0: <NO-CARRIER,BROADCAST,MULTICAST,UP> mtu 1500 qdisc mq state DOWN group default qlen 1000
    link/ether 06:8a:59:f4:f6:d3 brd ff:ff:ff:ff:ff:ff permaddr 24:ec:99:46:aa:1e
\end{lstlisting}
\end{latin}

مطابق خروجی فوق، سیستم مذکور دارای سه رابط شبکه مجزا است. نخستین مورد، رابط \cmd{lo} یا همان \en{Loopback} است؛ یک رابط مجازی که سیستم برای ارتباطات داخلی و محلی از آن بهره می‌برد. دومین مورد، \cmd{enp1s0}، رابط فیزیکی اترنت (\en{Ethernet}) و سومین مورد، \cmd{wlp2s0}، رابط شبکه‌ی بی‌سیم (\en{Wireless}) است.

در فرآیند عیب‌یابی (\en{Troubleshooting}) شبکه، دو شاخصه کلیدی باید مورد توجه قرار گیرند: نخست، عبارت \cmd{state UP} در خط اول هر رابط که نشان‌دهنده‌ی فعال بودن عملیاتی آن است؛ و دوم، وجود یک نشانی \en{IP} معتبر در مقابل فیلد \cmd{inet} در خط سوم. در سیستم‌هایی که از پروتکل \en{DHCP} جهت تخصیص خودکار آدرس استفاده می‌کنند، مشاهده‌ی یک نشانی \en{IP} معتبر در این بخش، گواهی بر عملکرد صحیح سرویس \en{DHCP} در شبکه است.

تحلیل جدول مسیریابی (\en{Routing Table}) گام بعدی در درک چگونگی هدایت ترافیک توسط سیستم‌عامل است. با استفاده از دستور زیر، می‌توانیم مسیرهای تعریف شده برای بسته‌های خروجی را مشاهده کنیم:

\begin{latin}
\begin{lstlisting}
[me@linuxbox ~]$ ip route show
default via 192.168.1.1 dev enp1s0 proto dhcp src 192.168.50.223 metric 100
169.254.0.0/16 dev enp1s0 scope link metric 1000
192.168.1.0/24 dev enp1s0 proto kernel scope link src 192.168.1.223 metric 100
\end{lstlisting}
\end{latin}

در خروجی نمونه، ساختار مسیریابی یک کلاینت در شبکه محلی (\en{LAN}) را مشاهده می‌کنیم. نشانی‌های \en{IP} که به صفر ختم می‌شوند (مانند \cmd{192.168.1.0/24})، نشان‌دهنده محدوده یک شبکه هستند. حضور این شبکه‌ها در جدول به این معناست که سیستم می‌تواند مستقیماً با تمامی میزبان‌های موجود در این بازه‌ها ارتباط برقرار کند.

در این خروجی، با دو شبکه مواجه هستیم: شبکه اصلی ما (\cmd{192.168.1.0}) و شبکه \cmd{169.254.0.0}. مورد دوم مربوط به مکانیزم \en{APIPA} (آدرس‌دهی خودکار خصوصی \en{IP}) است؛ این پروتکل زمانی فعال می‌شود که سرویس \en{DHCP} در شبکه در دسترس نباشد و سیستم به‌صورت خودکار یک نشانی موقت به خود اختصاص دهد تا ارتباطات پایه در لایه فیزیکی حفظ شود.

خط نخست که با واژه \cmd{default} آغاز می‌شود، حیاتی‌ترین بخش جدول است. این مسیر مشخص می‌کند که تمامی ترافیک‌هایی که مقصدشان در جدول تعریف نشده (مانند درخواست‌های اینترنتی)، باید به کدام سمت هدایت شوند. در اینجا، درگاه خروجی پیش‌فرض (\en{Default Gateway}) روی نشانی \cmd{192.168.1.1} (نشانی استاندارد مسیریاب‌های خانگی) تنظیم شده است که وظیفه هدایت بسته‌ها به شبکه گسترده‌تر را بر عهده دارد.

ابزار \cmd{ip} از نحو (\en{Syntax}) منعطفی پیروی می‌کند:

\begin{latin}
\begin{lstlisting}
ip [-options] object [command]
\end{lstlisting}
\end{latin}

در این ساختار، \file{object} می‌تواند مواردی نظیر \cmd{address} یا \cmd{route} باشد. یکی از ویژگی‌های کاربردی این ابزار، امکان خلاصه‌سازی دستورات است. از آنجا که دستور پیش‌فرض این برنامه \cmd{show} است، می‌توان عبارات را به شکل زیر کوتاه کرد:

\begin{itemize}
  \item \cmd{ip a}: معادل \cmd{ip address show}
  \item \cmd{ip r}: معادل \cmd{ip route show}
\end{itemize}

\section{انتقال فایل در بستر شبکه}

هدف غایی وجود یک شبکه، فراهم‌سازی امکان تبادل داده و جابجایی فایل‌ها میان گره‌های مختلف است. ابزارهای متعددی برای این امر توسعه یافته‌اند که در ادامه به بررسی یکی از قدیمی‌ترین آن‌ها پرداخته و در بخش‌های آتی، موارد پیشرفته‌تر را مرور خواهیم کرد.

\subsection{انتقال فایل با کلاینت ftp}

برنامه \cmd{ftp} یکی از راهکارهای کلاسیک دنیای شبکه است که نام خود را از پروتکل مورد استفاده‌اش، یعنی \en{File Transfer Protocol} (پروتکل انتقال فایل) وام گرفته است. پیش از ظهور و همه‌گیری مرورگرهای وب، \cmd{ftp} اصلی‌ترین روش دریافت فایل از اینترنت محسوب می‌شد؛ به‌طوری که بسیاری از مرورگرهای قدیمی‌تر از شناسه‌های \cmd{ftp://} پشتیبانی می‌کنند.

این ابزار جهت تعامل با سرویس‌دهندگان \en{FTP} (ماشین‌هایی که مخزن فایل‌های قابل بارگذاری و بارگیری هستند) طراحی شده است. لازم به ذکر است که پروتکل \en{FTP} در شکل اولیه‌ی خود فاقد امنیت است؛ چرا که اعتبارنامه‌ها (نام کاربری و رمز عبور) را به‌صورت متنی (\en{Cleartext}) و بدون رمزنگاری ارسال می‌کند. این نقیصه امنیتی باعث شده است که امروزه استفاده از آن در شبکه اینترنت، عمدتاً محدود به سرورهای \en{Anonymous} (ناشناس) باشد. در این سرورها، هر کاربری می‌تواند با نام کاربری \cmd{anonymous} و یک رمز عبور صوری وارد سیستم شود.

در نمونه‌ی زیر، یک نشست (\en{Session}) تعاملی جهت دریافت فایل ایمیج \en{ISO} توزیع اوبونتو از یک سرور فرضی نمایش داده شده است:

\begin{latin}
\begin{lstlisting}
[me@linuxbox ~]$ ftp fileservr
Connected to fileservr.localdomain.
220 (vsFTPD 2.0.1)
Name (fileservr:me): anonymous
331 Please specify the password.
Password:
230 Login successful.
Remote system type is UNIX.
Using binary mode to transfer files.
ftp> cd pub/cd_images/Ubuntu-24.04
250 Directory successfully changed.
ftp> ls
200 PORT command successful. Consider using PASV.
150 Here comes the directory listing.
-rw-rw-r--    1 500        500       733079552 Apr 25 03:53 ubuntu-24.04-desktop-amd64.iso
226 Directory send OK.
ftp> lcd Desktop
Local directory now /home/me/Desktop
ftp> get ubuntu-24.04-desktop-amd64.iso
local: ubuntu-24.04-desktop-amd64.iso remote: ubuntu-24.04-desktop-amd64.iso
200 PORT command successful. Consider using PASV.
150 Opening BINARY mode data connection for ubuntu-24.04-desktop-amd64.iso (733079552 bytes).
226 File send OK.
733079552 bytes received in 68.56 secs (10441.5 kB/s)
ftp> bye
\end{lstlisting}
\end{latin}

در جدول زیر، شرح عملکرد دستورات به‌کاررفته در نشست تعاملی فوق ارائه شده است:

\begin{table}[htbp]
\centering
\caption{دستورهای پرکاربرد در محیط تعاملی \cmd{ftp}}
\label{tab:ftp-commands}
\begin{tabularx}{\textwidth}{l X}
\toprule
\textbf{دستور} & \textbf{شرح عملکرد} \\
\midrule
\cmd{ftp fileserver} & فراخوانی کلاینت \cmd{ftp} و برقراریِ اتصال با سرویس‌دهنده‌ای به نام \en{fileserver}. در صورت موفقیت‌آمیز بودن اتصال، سرور بلافاصله درخواست نام کاربری می‌کند. \\
\addlinespace
\cmd{anonymous} & ورود به سیستم با استفاده از حساب کاربری ناشناس. پس از درج این نام، سرور درخواست گذرواژه می‌کند؛ برخی سرورها گذرواژه‌ی خالی را می‌پذیرند و برخی دیگر، نشانی ایمیل (مانند \cmd{user@example.com}) را به‌عنوان گذرواژه درخواست می‌کنند—هرچند این ایمیل معمولاً هیچ‌گاه اعتبارسنجی واقعی نمی‌شود. \\
\addlinespace
\cmd{cd pub/cd\_images/Ubuntu-24.04} & تغییرِ دایرکتوریِ جاری در سیستم راه دور (\en{Remote}). در اکثر سرورهای عمومی \en{FTP}، فایل‌های قابل دانلود برای عموم، بر اساس یک قرارداد رایج، در زیرشاخه‌ای به نام \cmd{pub} (مخفف \en{Public}) سازمان‌دهی می‌شوند. \\
\addlinespace
\cmd{ls} & مشاهده‌ی فهرست فایل‌ها و دایرکتوری‌های فرعی موجود در سیستم راه دور؛ مشابه عملکرد دستور \cmd{ls} در محیط پوسته، اما این‌بار برای مرور محتوای سرور. \\
\addlinespace
\cmd{lcd Desktop} & تغییر دایرکتوری جاری در سیستم محلی (\en{Local Change Directory}). برخلاف دستور \cmd{cd} که مسیر را روی سرور تغییر می‌دهد، \cmd{lcd} مسیر ذخیره‌سازی فایل‌ها را روی رایانه‌ی خود کاربر مشخص می‌کند؛ در این مثال، مسیر کاری کلاینت به دسکتاپ کاربر منتقل می‌شود تا فایل دریافتی در همان‌جا ذخیره گردد. \\
\addlinespace
\cmd{get ubuntu-24.04-desktop-amd64.iso} & ارسال درخواست انتقال یک فایل مشخص از سرور به سیستم محلی. فایل دریافتی، دقیقاً در دایرکتوری تعیین‌شده توسط آخرین دستور \cmd{lcd} ذخیره می‌شود؛ برای ارسال فایل در جهت عکس (از سیستم محلی به سرور)، از دستور متناظر \cmd{put} استفاده می‌شود. \\
\addlinespace
\cmd{bye} & خاتمه‌ی نشست و خروج از برنامه‌ی \cmd{ftp}، به‌همراه بستن اتصال با سرور. دستورهای \cmd{quit} و \cmd{exit} نیز عملکردی کاملاً مشابه دارند و هر سه به‌جای یکدیگر قابل استفاده‌اند. \\
\bottomrule
\end{tabularx}
\end{table}

با وارد کردن عبارت \cmd{help} در اعلانِ \cmd{ftp>} می‌توانید فهرست کاملی از دستورات پشتیبانی‌شده را مشاهده کنید. در صورتی که مجوزهای دسترسی (\en{Permissions}) لازم بر روی سرور صادر شده باشد، می‌توان بسیاری از عملیات‌های مدیریت فایل را از این طریق انجام داد؛ هرچند این شیوه در مقایسه با ابزارهای مدرن، تا حدی زمان‌بر و دشوار محسوب می‌شود.

\subsection{قابلیت‌ها و ویژگی‌های ابزار lftp}

برنامه \cmd{ftp} تنها کلاینت مبتنی بر خط فرمان برای تعامل با این پروتکل نیست؛ در واقع، گزینه‌های متعددی در این حوزه توسعه یافته‌اند که در میان آن‌ها، \cmd{lftp} (اثر الکساندر لوکیانوف) یکی از قدرتمندترین و محبوب‌ترین ابزارها به‌شمار می‌رود. این برنامه ضمن حفظ ساختار تعاملی مشابه با \cmd{ftp} سنتی، قابلیت‌های پیشرفته‌ای را ارائه می‌دهد که بهره‌وری و سهولت استفاده از آن را به‌مراتب افزایش داده است.

از جمله شاخص‌ترین ویژگی‌های این ابزار می‌توان به موارد زیر اشاره کرد:

\begin{itemize}
  \item \textbf{پشتیبانی از پروتکل‌های متعدد:} توانایی تعامل با پروتکل‌های مختلف از جمله \en{HTTP}، \en{SFTP} و \en{FISH}.
  \item \textbf{مدیریت هوشمند خطا و اتصال مجدد:} قابلیت تلاش مجدد خودکار (\en{Auto-retry}) در صورت قطع اتصال یا بروز اختلال در بارگیری.
  \item \textbf{چندوظیفگی و پردازش پس‌زمینه:} امکان اجرای فرآیندهای انتقال داده در پس‌زمینه (\en{Background Processing}).
  \item \textbf{تکمیل خودکار (\en{Tab Completion}):} پشتیبانی از کلید \en{Tab} برای تکمیل خودکار مسیرها و نام فایل‌ها، مشابه با محیط پوسته.
  \item \textbf{قابلیت \en{Mirroring}:} امکان همگام‌سازی کامل یک دایرکتوری محلی با دایرکتوری راه‌دور و برعکس (\en{Reverse Mirror}).
\end{itemize}

\subsection{انتقال داده از طریق URL با curl}

یکی دیگر از ابزارهای بسیار منعطف و پرکاربرد در حوزه‌ی انتقال فایل، \cmd{curl} است. در ساده‌ترین حالت استفاده، با تعیین یک نشانی اینترنتی (\en{URL})، این برنامه محتوای نخستین صفحه‌ی آن آدرس را دریافت کرده و در خروجی استاندارد (\en{Standard Output}) نمایش می‌دهد:

\begin{latin}
\begin{lstlisting}
[me@linuxbox ~]$ curl https://linuxcommand.org
\end{lstlisting}
\end{latin}

کلاینت \cmd{curl} علاوه بر توانایی دریافت هم‌زمان چندین \en{URL}، از طیف گسترده‌ای از پروتکل‌های شبکه از جمله \en{HTTP}، \en{HTTPS}، \en{FTP}، \en{IMAP}، \en{POP3}، \en{SFTP} و \en{SMB} پشتیبانی می‌کند. در جدول زیر، برخی از گزینه‌های پرکاربرد این ابزار آورده شده است:

\begin{table}[htbp]
\centering
\caption{گزینه‌های متداول در فرمان \cmd{curl}}
\label{tab:curl-options}
\begin{tabularx}{\textwidth}{l X}
\toprule
\textbf{گزینه} & \textbf{شرح عملکرد} \\
\midrule
\cmd{-o, --output file} & ذخیره‌ی خروجی در یک فایل مشخص به‌جای نمایش در ترمینال. \\
\addlinespace
\cmd{-O, --remote-name} & ذخیره‌ی فایل با نام اصلی آن در سرور راه دور. \\
\addlinespace
\cmd{-s, --silent} & اجرای دستور در حالت بی‌صدا (عدم نمایش نوار پیشرفت و خطاها). \\
\addlinespace
\cmd{-u, --proxy-user user:password} & تعیین اعتبارنامه‌های لازم (نام کاربری و رمز عبور) جهت احراز هویت. \\
\addlinespace
\cmd{-v, --verbose} & نمایش جزئیات کامل فرآیند ارتباط و هدرهای پروتکل (مناسب برای عیب‌یابی). \\
\bottomrule
\end{tabularx}
\end{table}

برای درک جزئیات پیچیده‌تر و آشنایی با تمامی قابلیت‌های این ابزار قدرتمند، مطالعه‌ی صفحه‌ی راهنمای آن (\cmd{man curl}) توصیه می‌شود.

\subsection{بارگیری غیرتعاملی فایل با wget}

یکی دیگر از ابزارهای قدرتمند خط فرمان جهت دریافت داده‌ها، \cmd{wget} است. این برنامه به‌طور تخصصی برای بارگیری محتوا از پروتکل‌های \en{HTTP}، \en{HTTPS} و \en{FTP} طراحی شده و قابلیت دریافت فایل‌های مجزا، مجموعه‌ای از فایل‌ها یا حتی کپی‌برداری کامل از وب‌سایت‌ها را داراست.

به‌عنوان مثال، برای بارگیری صفحه‌ی نخست وب‌سایت \file{linuxcommand.org} می‌توان از دستور زیر استفاده کرد:

\begin{latin}
\begin{lstlisting}
[me@linuxbox ~]$ wget http://linuxcommand.org/index.php
--11:02:51--  http://linuxcommand.org/index.php
           => `index.php'
Resolving linuxcommand.org... 66.35.250.210
Connecting to linuxcommand.org|66.35.250.210|:80... connected.
HTTP request sent, awaiting response... 200 OK
Length: unspecified [text/html]

[ <=>                                   ] 3,120         --.-K/s

11:02:51 (161.75 MB/s) - `index.php' saved [3120]
\end{lstlisting}
\end{latin}

تنوع گسترده‌ی گزینه‌های اجرایی در \cmd{wget} امکانات زیر را فراهم می‌آورد:

\begin{itemize}
  \item \textbf{بارگیری بازگشتی (\en{Recursive}):} جهت پیمایش خودکار پیوندها و دریافت سلسله‌مراتبی داده‌ها.
  \item \textbf{اجرا در پس‌زمینه (\en{Background}):} تداوم فرآیند بارگیری حتی پس از خروج کاربر از سیستم (\en{Logout}).
  \item \textbf{قابلیت \en{Resume}:} ازسرگیری خودکار فایل‌هایی که به‌دلیل اختلال در اتصال، به‌صورت ناقص بارگیری شده‌اند.
\end{itemize}

مستندات این ابزار در صفحه‌ی راهنمای آن (\cmd{man wget}) با دقتی فراتر از استانداردهای معمول تدوین شده و تمامی جزئیات فنی آن را پوشش می‌دهد.

\section{مدیریت و ارتباط امن با میزبان‌های راه دور}

سیستم‌عامل‌های خانواده یونیکس (\en{Unix-like}) از دهه‌های گذشته قابلیت مدیریت از راه دور از طریق شبکه را در بطن خود داشته‌اند. در دوران آغازین توسعه‌ی شبکه و پیش از همه‌گیری اینترنت، ابزارهایی نظیر \cmd{rlogin} و \cmd{telnet} از محبوب‌ترین راهکارهای ورود به سیستم‌های راه دور (\en{Remote Login}) محسوب می‌شدند. با این حال، این پروتکل‌ها از همان حفره‌ی امنیتی ابزار \cmd{ftp} رنج می‌برند؛ تمامی داده‌های تبادل شده، از جمله نام‌های کاربری و رمزهای عبور، به‌صورت متنی و فاقد رمزنگاری (\en{Plaintext}) در شبکه منتقل می‌شوند. این ضعف ساختاری، استفاده از آن‌ها را در فضای ناامن اینترنت مدرن کاملاً منسوخ و غیرقابل‌قبول کرده است.

\subsection{معماری و نحوه عملکرد پروتکل SSH}

برای مرتفع ساختن چالش‌های امنیتی، پروتکل \en{SSH} (مخفف \en{Secure Shell}) توسعه یافت. این پروتکل دو معضل بنیادین در ارتباطات راه دور را به طور کامل حل می‌کند:

\begin{enumerate}
  \item \textbf{تأیید اصالت میزبان:} اطمینان از هویت واقعی سرور مقصد که مانع از وقوع حملات «مرد میانی» (\en{Man-in-the-Middle}) می‌شود.
  \item \textbf{رمزنگاری سراسری:} تمامی داده‌های تبادل شده میان میزبان محلی و راه دور به‌صورت رمزنگاری‌شده منتقل می‌شوند.
\end{enumerate}

ساختار \en{SSH} مبتنی بر مدل کلاینت-سرور است. یک سرویس‌دهنده (\en{SSH Server}) روی میزبان راه دور اجرا شده و به‌صورت پیش‌فرض بر روی پورت ۲۲ به درخواست‌ها پاسخ می‌دهد؛ در سوی دیگر، کاربر از یک برنامه کلاینت (\en{SSH Client}) برای برقراری اتصال استفاده می‌کند.

بسیاری از توزیع‌های لینوکس از نسخه‌ی \en{OpenSSH} بهره می‌برند که پروژه‌ای تحت نظر بنیاد \en{OpenBSD} است. برخی توزیع‌ها هر دو بسته‌ی کلاینت و سرور را به‌صورت پیش‌فرض ارائه می‌دهند، اما در برخی دیگر تنها کلاینت نصب شده است. برای آنکه یک سیستم قابلیت پذیرش اتصالات راه دور را داشته باشد، باید بسته‌ی \cmd{openssh-server} بر روی آن نصب و فعال گردد؛ همچنین در صورت فعال بودن دیوار آتش (\en{Firewall})، باید اجازه‌ی عبور ترافیک از پورت \en{TCP/22} صادر شود.

\begin{note}
\textbf{نکته:} چنانچه سیستم راه دوری برای آزمایش در اختیار ندارید، می‌توانید با نصب بسته‌ی سرور روی ماشین خود، از کلمه‌ی کلیدی \cmd{localhost} به عنوان نام میزبان استفاده کنید. در این حالت، سیستم یک اتصال شبکه‌ای داخلی با خودش برقرار می‌کند.
\end{note}

برنامه کلاینت که وظیفه برقراری ارتباط با سرویس‌دهنده‌های راه دور را بر عهده دارد، به سادگی \cmd{ssh} نامیده می‌شود. برای اتصال به میزبانی با نام فرضی \file{remote-sys} از دستور زیر استفاده می‌کنیم:

\begin{latin}
\begin{lstlisting}
[me@linuxbox ~]$ ssh remote-sys
The authenticity of host 'remote-sys (192.168.1.4)' can't be established.
RSA key fingerprint is
41:ed:7a:df:23:19:bf:3c:a5:17:bc:61:b3:7f:d9:bb.
Are you sure you want to continue connecting (yes/no)?
\end{lstlisting}
\end{latin}

در اولین تلاش برای برقراری اتصال، پیامی مبنی بر عدم امکان تایید اصالت میزبان نمایش داده می‌شود؛ چرا که کلاینت پیش‌تر با این سرور تعاملی نداشته و شناسه‌ای از آن در اختیار ندارد. جهت پذیرش اعتبارنامه سرور، باید عبارت \cmd{yes} را تایپ کنید. پس از تایید، سیستم از شما گذرواژه می‌خواهد:

\begin{latin}
\begin{lstlisting}
Warning: Permanently added 'remote-sys,192.168.1.4' (RSA) to the list of known hosts.
me@remote-sys's password:
\end{lstlisting}
\end{latin}

پس از درج صحیح گذرواژه، اعلان خط فرمان (\en{Prompt}) مربوط به پوسته‌ی سیستم راه دور ظاهر می‌گردد:

\begin{latin}
\begin{lstlisting}
Last login: Sat Aug 30 13:00:48 2024
[me@remote-sys ~]$
\end{lstlisting}
\end{latin}

نشست پوسته راه دور تا زمانی که کاربر دستور \cmd{exit} را وارد نکند، فعال باقی می‌ماند. با بستن اتصال، نشست پوسته محلی مجدداً از سر گرفته شده و اعلان خط فرمان سیستم خودتان دوباره نمایان می‌شود.

همچنین، امکان اتصال به سیستم‌های راه دور با نام کاربری متفاوت از نام محلی نیز فراهم است. به‌طور مثال، اگر کاربری قصد ورود به حساب کاربری دیگری در سیستم میزبان مقصد را داشته باشد، می‌تواند از ساختار زیر استفاده کند:

\begin{latin}
\begin{lstlisting}
[me@linuxbox ~]$ ssh bob@remote-sys
bob@remote-sys's password:
Last login: Sat Aug 30 13:03:21 2024
[bob@remote-sys ~]$
\end{lstlisting}
\end{latin}

همان‌طور که پیش‌تر اشاره شد، پروتکل \en{SSH} وظیفه‌ی تأیید اصالت میزبان راه دور را بر عهده دارد. چنانچه میزبان مقصد نتواند هویت خود را با کلیدهای ثبت‌شده‌ی قبلی تطبیق دهد، هشدار امنیتی زیر نمایش داده می‌شود:

\begin{latin}
\begin{lstlisting}
[me@linuxbox ~]$ ssh remote-sys
@@@@@@@@@@@@@@@@@@@@@@@@@@@@@@@@@@@@@@@@@@@@@@@@@@@@@@@@@@@
@    WARNING: REMOTE HOST IDENTIFICATION HAS CHANGED!     @
@@@@@@@@@@@@@@@@@@@@@@@@@@@@@@@@@@@@@@@@@@@@@@@@@@@@@@@@@@@
IT IS POSSIBLE THAT SOMEONE IS DOING SOMETHING NASTY!
Someone could be eavesdropping on you right now (man-in-the-middle attack)!
It is also possible that a host key has just been changed.
The fingerprint for the ECDSA key sent by the remote host is
SHA256:9uhOXHpZp0AjIs0ZDAnNYfv7ksU04Mjy5pR4kUTSKKA.
Please contact your system administrator.
Add correct host key in /home/me/.ssh/known_hosts to get rid of this message.
Offending ECDSA key in /home/me/.ssh/known_hosts:42
 remove with:
  ssh-keygen -f "/home/me/.ssh/known_hosts" -R "remote-sys"
Host key for remote-sys has changed and you have requested strict checking.
Host key verification failed.
\end{lstlisting}
\end{latin}

این پیام هشداردهنده نشان‌دهنده‌ی وقوع یکی از دو سناریوی محتمل است: نخست، احتمال وقوع یک حمله‌ی «مرد میانی» (\en{Man-in-the-Middle}) وجود دارد که در آن مهاجم سعی می‌کند ترافیک شما را شنود یا منحرف کند. هرچند به دلیل هوشیاری ابزار \en{SSH} در آگاه‌سازی کاربر، این نوع حملات کمتر رخ می‌دهند. سناریوی دوم که محتمل‌تر به نظر می‌رسد، بروز تغییرات ساختاری در میزبان راه دور است؛ برای مثال ممکن است سیستم‌عامل سرور مجدداً نصب شده یا تنظیمات سرویس‌دهنده‌ی \en{SSH} تغییر یافته باشد که منجر به تولید کلید شناسایی (\en{Host Key}) جدید شده است.

با این وجود، به دلیل حساسیت‌های امنیتی، هرگز نباید احتمال وقوع حمله را به‌سادگی نادیده گرفت. در صورت مواجهه با این وضعیت، الزامی است پیش از هر اقدامی با مدیر سیستمِ مقصد تماس گرفته و از صحت تغییرات اطمینان حاصل کنید.

پس از اطمینان از اینکه منشأ پیام هشدار فاقد مخاطرات امنیتی است، می‌توان اختلال پیش‌آمده را از سمت کلاینت مرتفع نمود. ساده‌ترین روش، استفاده از دستوری است که در متن هشدار پیشنهاد شده است:

\begin{latin}
\begin{lstlisting}
ssh-keygen -f "/home/me/.ssh/known_hosts" -R "remote-sys"
\end{lstlisting}
\end{latin}

در غیر این صورت، می‌توان به‌صورت دستی و با استفاده از یک ویرایشگر متن (نظیر \cmd{vim})، کلید منسوخ را از فایل \file{\textasciitilde/.ssh/known\_hosts} حذف کرد. در پیام هشدار نمونه، عبارت زیر راهنمای ما خواهد بود:

\begin{latin}
\begin{lstlisting}
Offending key in /home/me/.ssh/known_hosts:42
\end{lstlisting}
\end{latin}

این عبارت بیانگر آن است که کلید ناسازگار در خط ۴۲ فایل \file{known\_hosts} قرار دارد. با حذف این خط، برنامه \cmd{ssh} قادر خواهد بود در اتصال بعدی، اعتبارنامه‌های جدید میزبان راه دور را دریافت و ذخیره کند.

علاوه بر امکان برقراری یک نشست تعاملی کامل، ابزار \cmd{ssh} قابلیت اجرای یک «تک‌فرمان» بر روی سیستم مقصد را نیز داراست. به‌طور مثال، جهت اجرای دستور \cmd{free} در میزبان \file{remote-sys} و مشاهده‌ی وضعیت حافظه آن در ترمینال محلی، دستور زیر به‌کار می‌رود:

\begin{latin}
\begin{lstlisting}
[me@linuxbox ~]$ ssh remote-sys free
me@twin4's password:
        total         used        free      shared    buffers     cached
Mem:   775536       507184      268352           0     110068     154596
-/+ buffers/cache:  242520      533016
Swap: 1572856            0     1572856
[me@linuxbox ~]$
\end{lstlisting}
\end{latin}

این تکنیک را می‌توان برای سناریوهای کاربردی‌تری نیز به خدمت گرفت؛ مانند مثال زیر که در آن دستور \cmd{ls} روی سیستم راه دور اجرا شده و خروجی آن به فایلی در سیستم محلی هدایت (\en{Redirect}) می‌شود:

\begin{latin}
\begin{lstlisting}
[me@linuxbox ~]$ ssh remote-sys 'ls *' > dirlist.txt
me@twin4's password:
[me@linuxbox ~]$
\end{lstlisting}
\end{latin}

\begin{note}
\textbf{نکته حائز اهمیت:} استفاده از تک‌کوتیشن‌ها (\cmd{' '}) پیرامون فرمان راه دور است. این کار مانع از تفسیر و «بسط مسیر» (\en{Pathname Expansion}) توسط پوسته‌ی سیستم محلی شده و اجازه می‌دهد تا الگوها (مانند \cmd{*}) مستقیماً در سیستم مقصد پردازش شوند. همچنین، چنانچه هدف ما ذخیره‌ی خروجی در فایلی بر روی خودِ ماشین راه دور باشد، باید عملگر ریدایرکت را نیز به همراه نام فایل درون کوتیشن قرار دهیم:
\begin{latin}
\begin{lstlisting}
[me@linuxbox ~]$ ssh remote-sys 'ls * > dirlist.txt'
\end{lstlisting}
\end{latin}
\end{note}

\subsection{تونل‌زنی و هدایت ترافیک در SSH}

فرآیند برقراری اتصال با یک میزبان راه دور از طریق \en{SSH}، با ایجاد یک تونل رمزنگاری‌شده (\en{Encrypted Tunnel}) میان سیستم محلی و مقصد همراه است. در حالت معمول، این تونل وظیفه‌ی انتقال امن فرامین صادر شده از مبدأ به مقصد و بازگرداندن نتایج به سیستم محلی را بر عهده دارد. افزون بر این کاربرد اولیه، پروتکل \en{SSH} این قابلیت را دارد که انواع ترافیک شبکه را از میان این تونل عبور دهد و عملاً نوعی شبکه خصوصی مجازی (\en{VPN}) موقت میان دو گره ایجاد نماید.

یکی از رایج‌ترین کاربردهای این قابلیت، انتقال ترافیک گرافیکی تحت سیستم پنجره‌بندی \en{X} (موسوم به \en{X Window System}) است. در سیستم‌هایی که همچنان از \en{X Server} بهره می‌برند (برخلاف توزیع‌های مدرن مبتنی بر \en{Wayland})، می‌توان یک برنامه کلاینت \en{X} (یک نرم‌افزار با واسط گرافیکی) را روی سرور راه دور اجرا کرد و خروجی بصری آن را در دسکتاپ محلی مشاهده نمود. برای درک بهتر، فرض کنید پشت سیستمی به نام \file{linuxbox} نشسته‌اید و قصد دارید ابزار گرافیکی \cmd{xload} را از روی سرور \file{remote-sys} فراخوانی کرده و پنجره‌ی آن را در مانیتور خود ببینید. این عملیات با استفاده از گزینه \cmd{-X} به سادگی انجام می‌شود:

\begin{latin}
\begin{lstlisting}
[me@linuxbox ~]$ ssh -X remote-sys
me@remote-sys's password:
Last login: Mon Sep 08 13:23:11 2024
[me@remote-sys ~]$ xload
\end{lstlisting}
\end{latin}

به محض اجرای فرمان \cmd{xload} در ترمینال راه دور، پنجره‌ی گرافیکی آن در محیط دسکتاپ محلی شما ظاهر خواهد شد. شایان ذکر است که در برخی پیکربندی‌های خاص سیستم، ممکن است به‌جای گزینه \cmd{-X} نیاز به استفاده از گزینه \cmd{-Y} (جهت برقراری اتصال با سطوح دسترسی قابل‌اعتمادتر یا \en{Trusted X11 Forwarding}) باشد.

\subsection{انتقال امن فایل با ابزارهای scp و sftp}

بسته‌ی \en{OpenSSH} علاوه بر فراهم‌سازی دسترسی راه دور، شامل دو ابزار کاربردی است که از تونل رمزنگاری‌شده‌ی \en{SSH} برای جابجایی فایل‌ها بهره می‌برند. نخستین ابزار، \cmd{scp} (مخفف \en{Secure Copy}) است که عملکردی مشابه دستور استاندارد \cmd{cp} دارد؛ با این تفاوت کلیدی که مسیر مبدأ یا مقصد می‌تواند در یک سیستم راه دور باشد. در این ساختار، آدرس‌دهی با نام میزبان آغاز شده و با علامت دو نقطه (\cmd{:}) از مسیر فایل جدا می‌شود. به‌طور مثال، برای کپی کردن فایل \file{document.txt} از دایرکتوری خانگی در سرور \file{remote-sys} به دایرکتوری جاری در سیستم محلی، دستور زیر به‌کار می‌رود:

\begin{latin}
\begin{lstlisting}
[me@linuxbox ~]$ scp remote-sys:document.txt .
me@remote-sys's password:
document.txt                       100%  5581     5.5KB/s   00:00
[me@linuxbox ~]$
\end{lstlisting}
\end{latin}

مشابه فرمان \cmd{ssh}، چنانچه نام کاربری شما در سیستم مقصد متفاوت است، می‌توانید آن را پیش از نام میزبان درج کنید:

\begin{latin}
\begin{lstlisting}
[me@linuxbox ~]$ scp bob@remote-sys:document.txt .
\end{lstlisting}
\end{latin}

دومین ابزار، \cmd{sftp} است که همان‌طور که از نامش پیداست، به‌عنوان جایگزینی امن برای پروتکل کلاسیک \cmd{ftp} توسعه یافته است. ابزار \cmd{sftp} تمامی عملیات را در بستر تونل \en{SSH} انجام می‌دهد و مزیت راهبردی آن نسبت به \cmd{ftp} این است که به هیچ سرویس‌دهنده‌ی مجزایی نیاز ندارد؛ در واقع هر سیستمی که از \en{SSH Server} بهره می‌برد، به‌طور خودکار قابلیت تبادل فایل از طریق \cmd{sftp} را نیز خواهد داشت. در ادامه، نمونه‌ای از یک نشست تعاملی \cmd{sftp} آورده شده است:

\begin{latin}
\begin{lstlisting}
[me@linuxbox ~]$ sftp remote-sys
Connecting to remote-sys...
me@remote-sys's password:
sftp> ls
ubuntu-8.04-desktop-i386.iso
sftp> lcd Desktop
sftp> get ubuntu-8.04-desktop-i386.iso
Fetching /home/me/ubuntu-8.04-desktop-i386.iso to ubuntu-8.04-desktop-i386.iso
/home/me/ubuntu-8.04-desktop-i386.iso 100% 699MB   7.4MB/s   01:35
sftp> bye
\end{lstlisting}
\end{latin}

\begin{note}
\textbf{نکته:} پروتکل \en{SFTP} به‌طور گسترده توسط مدیران فایل گرافیکی (\en{File Managers}) در لینوکس پشتیبانی می‌شود. در محیط‌هایی نظیر \en{GNOME} (\en{Nautilus}) یا \en{KDE} (\en{Dolphin})، می‌توانید با وارد کردن پروتکل \cmd{sftp://} در نوار آدرس (مانند \cmd{sftp://remote-sys})، به فایل‌های موجود در سرور راه دور به‌صورت بصری و مشابه درایوهای محلی دسترسی داشته باشید.
\end{note}

\section{کلاینت‌های SSH در سیستم‌عامل ویندوز}

چنانچه در محیط ویندوز نیاز به دسترسی و مدیریت یک سرور لینوکسی داشته باشید، گام نخست تهیه و استفاده از یک برنامه کلاینت \en{SSH} متناسب با این سیستم‌عامل است. در میان ابزارهای متعدد توسعه یافته، \cmd{PuTTY} (اثر سایمون تاتهام و تیم وی) به‌عنوان محبوب‌ترین و شناخته‌شده‌ترین گزینه در این حوزه به شمار می‌رود. \cmd{PuTTY} با شبیه‌سازی یک محیط ترمینال، امکان برقراری نشست‌های \en{SSH} (و همچنین \en{Telnet}) را برای کاربر فراهم می‌کند. علاوه بر این، این مجموعه ابزارهای جانبی برای پروتکل‌های \cmd{scp} و \cmd{sftp} را نیز جهت انتقال امن فایل در اختیار کاربر قرار می‌دهد.

علاقه‌مندان می‌توانند این ابزار را از نشانی زیر دریافت کنند:

\begin{latin}
\url{http://www.chiark.greenend.org.uk/~sgtatham/putty/}
\end{latin}

\section{جمع‌بندی}

در این فصل، به بررسی کلیدی‌ترین ابزارهای شبکه که به‌صورت پیش‌فرض در اغلب توزیع‌های لینوکسی در دسترس هستند، پرداختیم. با توجه به جایگاه محوری لینوکس در زیرساخت‌های سرور و تجهیزات شبکه، ابزارهای تخصصی بی‌شمار دیگری نیز وجود دارند که در صورت نیاز قابل نصب و بهره‌برداری هستند. با این حال، تسلط بر همین مجموعه ابزارهای پایه، توانایی لازم برای انجام طیف گسترده‌ای از وظایف مدیریتی و عیب‌یابی شبکه را برای کاربر فراهم می‌آورد.

\section{منابع مطالعه تکمیلی}

برای یک دید کلی (هرچند قدیمی) به مبحث مدیریت شبکه، پروژه مستندات لینوکس (\en{LDP}) راهنمای \en{``Linux Network Administrator's Guide''} را ارائه می‌کند:

\begin{latin}
\url{http://tldp.org/LDP/nag2/index.html}
\end{latin}

ویکی‌پدیا مقالات متعددی در زمینه شبکه‌های کامپیوتری ارائه می‌دهد. برخی از موارد پایه عبارتند از:

\begin{latin}
\url{http://en.wikipedia.org/wiki/Internet_protocol_address}\\
\url{http://en.wikipedia.org/wiki/Host_name}\\
\url{http://en.wikipedia.org/wiki/Uniform_Resource_Identifier}
\end{latin}

صفحه اصلی پروژه \en{curl} یک آموزش عالی ارائه می‌دهد که نمونه‌های کاربردی فراوانی از قابلیت‌های این ابزار را نمایش می‌دهد:

\begin{latin}
\url{https://curl.se/docs/tutorial.html}
\end{latin}

برای درک چگونگی و چرایی آدرس‌دهی خودکار \en{IP} خصوصی (\en{APIPA})، پیوند زیر توضیحات مناسبی ارائه می‌دهد:

\begin{latin}
\url{https://www.33rdsquare.com/what-does-169-254-0-0-mean/}
\end{latin}